
\section{Schéma Monolithique (DG/FV)}
La méthode se repose sur deux détails clés : une réécriture du schéma DG utilisant des sous-mailles et 
un flux hybride entre le flux DG du schéma réécrit et un flux Volume Fini (VF) entre sous-mailles.
\subsection{Réecriture du schéma DG comme schéma sur sous-mailles.}
\paragraph{sous-maillage \\}

On dote chaque maille $\omega_i$ d'un sous-maillage : $\{S_m^i\}_m = ]\tilde{x}^i_{m-1/2}, \tilde{x}^i_{m+1/2}[$ vérifiant les mêmes propriétés qu'un maillage \eqref{maillage} où $\bigcup_m S_m^i = \omega_i$ et $\tilde{x}_{1/2}^i = x_{i-1/2},\tilde{x}_{N_s^i + 1/2}^i = x_{i+1/2} $\\
On cherchera à passer les degrées de libertés des polynomes de $ \sigma \in\mathbb{P}^k(\omega_i)$, des moments sur la base polynomiale $ \left\{ \underline{\sigma}_p \right\}_p$ aux moyennes de ces polynomes sur le sous-mailles $\{S_m^i \}$, $ \left\{\overline{\sigma}^m \right\}_m$ où $\overline{\sigma}^m := \int_{\omega_i} {\sigma}^m(x) dx$  \\
Soit la matrice $$\mathcal{P}_{m,p} = \frac{1}{|S_m^i|} \int_{S_m^i} \sigma_p(x) dx$$ 
Pour que le passage des moments polynomiaux $\left\{ \underline{\sigma}_p \right\}_p$ aux sous-volumes finis soit injectif, il nous faut : $N_s \geq N_k$. 
De manière générale, on dit que le sous-maillage est \textit{admissible} si et seulement si la matrice $\mathcal{P}^t\,\mathcal{P}$ est inversible.

\paragraph{Remarque \\}  
\begin{wrapfigure}{r}{0.4\textwidth} 
    \includegraphics[width =0.3\textwidth]{schema/projection_L2_svf.png}
    \caption{Comparaison entre $\underline{U}$ et $\overline{U}$.}
\end{wrapfigure}
Pour avoir une équivalence au sens le plus strict possible, il nous faut $N_s = N_k$. 
Dans le cas échéant ($N_s > N_k$), on utilisera un opérateur de reconstruction $\mathcal{R} = \left(\mathcal{P}^t\,\mathcal{P}\right)^{-1} \mathcal{P}^t$  qui nous donnera les moments polynomiaux au sens des moindres carrés. \\
Il est alors simple de passer de la représentation polynomiale $\underline{U}$ à la représentation en sous-moyennes $\overline{U}$ ou inversement par ces relations: 
\begin{equation}\label{passage_svf}
    \overline{U} = \mathcal{P}\underline{U} \qquad \underline{U} = \mathcal{R}\overline{U}
\end{equation}
On en profite pour introduire les projections sur $\mathbb{P}^k(\omega_i)$ de $\{\mathbb{1}_{S^i_m}\}_{1\leq m\leq N_s}$ noté $\{\phi_m^i\}_{1\leq m\leq N_s}$ 
$$\int_{\omega_i} \phi_m^i \sigma dx = \int_{S_m^i} \sigma dx, \, \forall \sigma \in \mathbb{P}^k(\omega_i)$$
les intégrales sur $S_m^i$ seront calculées avec une méthode de quadrature sur $S_m^i$ et non sur $\omega_i$.

\paragraph{initialisation de $\overline{U}$}
Il est important de noter que contrairement à la méthode DG "classique", il est possible et grandement utile de changer la méthode d'initialisation.\\
En effet, si l'on prend naivement les sous-moyennes des polynômes que l'on a calculé comme projection $\mathbb{L}^2$ de $u_0$, on peut avoir $\overline{U} \notin \Omega$.
Ainsi, il préférable d'initialiser $\overline{U}$ comme sous-moyennes de $U_0$ et de définir $\underline{U}$ à partir de $\overline{U}$ par moindres carrés \eqref{passage_svf} 
\paragraph{ Construction des flux reconstruit \red{nom à changer} \\}

En reprenant la formulation variationnelle \eqref{Formulation_Variationnelle}, si l'on projette ${f}({u_h})$ sur $\mathbb{P}^k(\omega_i)$ : ${F}_h$, on obtient une formulation équivalente : 
$$\int_{\omega_i} \partial_t {u_h} \sigma dx - \int_{\omega_i} {F}_h^i \partial_x \sigma dx 
+ \left[ \mathcal{F}\, \sigma \right]_{x_{i-\frac{1}{2}}}^{x_{i+\frac{1}{2}}} = 0 \qquad \forall  \sigma \in \mathbb{P}^k(\omega_i) $$  
Alors par le biais d'une intégration par partie, on obtient la forme \textit{forte} des schéma DG :  
$$\int_{\omega_i} \partial_t {u_h} \sigma dx = \int_{\omega_i} \partial_x {F}_h^i  \sigma dx 
- \left[ ( {F}_h^i - \mathcal{F} )\, \sigma \right]_{x_{i-\frac{1}{2}}}^{x_{i+\frac{1}{2}}} = 0 \qquad \forall  \sigma \in \mathbb{P}^k(\omega_i) $$  

Comme cette équation est vrai pour tout polynôme de $\mathbb{P}^k(\omega_i)$, il est en particulier vrai pour $\phi_m^i$. Cela nous amène presque à un flux ne dépendant que de $S_m^i$ : 
\begin{equation}
    \int_{S_m^i} \partial_t {u_h} dx = - \left(\int_{S_m^i}  \partial_x {F}_h^i dx - \left[ ( {F}_h^i - \mathcal{F} )\, \phi_m \right]_{x_{i-\frac{1}{2}}}^{x_{i+\frac{1}{2}}}\right) \qquad \forall m\in \{ 1, \dots, N_s\}
\end{equation}
Il est possible d'obtenir un schéma volume fini sur les $\{\overline{{u}_h}^m\}$ en fabriquant un flux avec pour unique but de ressembler à un flux volume fini. \\
Soit $\widehat{{F}^i}_{m+1/2}$ ces flux reconstruits posés tels que : 

\begin{align}
&\widehat{{F}^i}_{m+1/2} - \widehat{{F}^i}_{m-1/2} =  -  \left[{F}_h^i\right]_{\tilde{x}^i_{m-1/2}}^{\tilde{x}^i_{m+1/2}}  - \left[ ( {F}_h^i - \mathcal{F} )\, \phi_m \right]_{x_{i-\frac{1}{2}}}^{x_{i+\frac{1}{2}}} \qquad \forall m\in \{ 1, \dots, N_s\} \\
& \text{où l'on impose : } \widehat{{F}^i}_{1/2} = \mathcal{F}_{i-1/2} \qquad \widehat{{F}^i}_{ N_s + 1/2} = \mathcal{F}_{i+1/2}
\end{align}
Il est possible d'obtenir une forme close de $\widehat{F^i}_{m+1/2}$: 
\begin{align}
    &\widehat{{F}^i}_{m+1/2} = {F}_h^i(\tilde{x}^i_{m+1/2}) 
    -C^{i-1/2}_{m+1/2} ( {F}_h^i(x_{i-1/2}) - \mathcal{F}_{i-1/2}) 
    -C^{i+1/2}_{m+1/2} ( {F}_h^i(x_{i+1/2}) - \mathcal{F}_{i+1/2}) \\
    & \text{où }  C^{i-1/2}_{m+1/2} = (1- \sum_{p=1}^{m}\phi_p^i(x_{i-1/2})) \, , \qquad
        C^{i+1/2}_{m+1/2} =     \sum_{p=1}^{m}\phi_p^i(x_{i+1/2})
\end{align}

Tout ceci nous permet d'écrire le schéma DG sous une forme VF : 
 
\begin{equation}\label{DG_as_FV}
    \partial_t \overline{{u}_m}^i = -\frac{1}{|S_m^i|}(\widehat{{F}^i}_{m+1/2} - \widehat{{F}^i}_{m-1/2})
\end{equation}

\paragraph{Remarque \\}
\noindent Même si la forme de \eqref{DG_as_FV} laisse à penser que l'évolution de $\overline{{u_h}}^m$ 
ne dépend uniquement des sous-cellules voisines à $S_m^i$, 
il faut rappeler que $\widehat{{F}^i}_{m+1/2}$ dépend non seulement des sous-cellules voisines, 
mais aussi des cellules voisines $\omega_{i\pm1}$ par le biais du flux numérique évalué aux interfaces des cellules $x_{i\pm1/2}$ \\
Les flux $\widehat{{F}^i}_{m+1/2}$ pouvant se trouver entre deux sous-cellules appartenant à la même maille ou deux mailles différentes, 
Par soucis de claireté, nous écrirons $\widehat{{F}_{mp}}$ le flux allant de la sous-cellule $S_m^i$ à la sous-cellule $S_p^j$, $j\in \{i-1,i,i+1\}$ 

\subsection{Flux hybride (DG/FV)}
\noindent
Maintenant que nous avons établi un schéma agissant sur les sous valeurs moyennes
équivalent au schéma DG agissant sur les moments polynomiaux. 
Il devient alors possible de combiner les flux DG reconstruit $\widehat{F_{mp}}$ 
avec des flux proprement volumes finis défini sur les sous-mailles $F^{\text{FV}}_{mp}$ sous la forme d'un flux dit \textit{hybride} $\widetilde{F_{mp}}$ comme combinaison convexe des deux flux.

\begin{equation} 
\begin{aligned}
    \widetilde{F_{mp}}  &= F^{\text{FV}}_{mp} + \theta_{mp}\,\, (\widehat{F_{mp}} - F^{\text{FV}}_{mp})\\
                        &= F^{\text{FV}}_{mp} + \theta_{mp}\,\, \Delta F_{mp}
\end{aligned}
\end{equation}
Le schéma dépend alors du choix de $\theta_{mp}$. \\
Si l'on pose $\theta_{mp} = 1$, nous obtiendrons un schéma DG qui sera d'ordre élevé mais qui n'aura pas les propriétés souhaitées. \\
Si l'on pose $\theta_{mp} = 0$, nous obtiendrons un schéma VF qui sera entropique et suivra le principe du maximum local mais il sera d'ordre 1. \\
Avant de s'attarder sur le choix des $\theta_{mp}$, regardons le schéma que nous avons obtenus : 

\begin{equation}\label{DG_as_FV}
    \partial_t \overline{{u}_m}^i =  - \frac{1}{|S_m^i|} \sum_{S_p^j \in \mathcal{V}_m^i} \widetilde{F_{mp}}
\end{equation}

En se souvenant que toutes propriétés d'un schéma \textit{RK-SSP} d'ordre élevé est dérivée des propriétés du schéma d'Euler explicite, 
On peut voir que $\overline{{u}_m}^{i,n+1}$ peut s'écrire comme combinaison convexe entre $\overline{{u}_m}^{i,n}$ et ce que l'on définira comme des états de Riemann intermédiaires hybrides : 
$ \widetilde{{u}}_{mp}^- = \overline{{u}_m}^{i,n} - \frac{\widetilde{F_{mp}} -  F_h^i(\overline{{u}_m}^{i,n})\cdot n_{mp} }{\gamma_{mp}}$

\begin{align*}
&\overline{{u}_m}^{i,n+1} = \overline{{u}_m}^{i,n} - \frac{\Delta t_n}{|S_m^i|} \sum_{S_p^j \in \mathcal{V}_m^i} \widetilde{F_{mp}} 
\pm \frac{\Delta t_n}{|S_m^i|} \sum_{S_p^j \in \mathcal{V}_m^i} \gamma_{mp}\overline{{u}_m}^{i,n} \pm \frac{\Delta t_n}{|S_m^i|} F_h^i(\overline{{u}_m}^{i,n}) \\
&\overline{{u}_m}^{i,n+1}= \left( 1-  \frac{\Delta t_n}{|S_m^i|} \sum_{S_p^j \in \mathcal{V}_m^i} \gamma_{mp}\right) \overline{{u}_m}^{i,n} 
                                + \frac{\Delta t_n}{|S_m^i|} \sum_{S_p^j \in \mathcal{V}_m^i} \gamma_{mp} \left( \overline{{u}_m}^{i,n} - \frac{\widetilde{F_{mp}} -  F_h^i(\overline{{u}_m}^{i,n})\cdot n_{mp} }{\gamma_{mp}} \right)
\end{align*}
On note ici que la combinaison est convexe seulement si $ \frac{\Delta t_n}{|S_m^i|} \sum_{S_p^j \in \mathcal{V}_m^i} \gamma_{mp} \leq 1$, ce qui nous donne une condition CFL : 

\begin{equation}
    \Delta t \leq \frac{|S_m^i|}{\sum_{S_p^j \in \mathcal{V}_m^i} \gamma_{mp}}
\end{equation} 
Elle sera par ailleurs la seule condition CFL à prendre en compte car plus stricte que la condition CFL du schéma DG.\\ 

Il est souvent plus simple d'étudier les combinaisons convexes de $\overline{{u}_m}^{i,n+1}$ en découpant $ \widetilde{{u}}_{mp}^-$ comme : 
\begin{equation}
\begin{aligned}
    \widetilde{{u}}_{mp}^- =& \left( \frac{\overline{{u}_m}^{i,n}+ \overline{{u}_p}^{j,n}}{2} - \frac{(F_h^i(\overline{{u}_p}^{j,n}) - F_h^i(\overline{{u}_m}^{i,n})   )\cdot n_{mp}}{2 \gamma_{mp}}   \right) - \theta_{mp} \frac{\Delta F_{mp}}{\gamma_{mp}} \\
    =& {u}_{mp}^{*,FV} - \theta_{mp} \frac{\Delta F_{mp}}{\gamma_{mp}}
\end{aligned}    
\end{equation}

On note au passage ces égalité de symmétrie et d'antisymmétrie : \\
$n_{mp}= -n_{pm}$, et donc : $F^{\text{FV}}_{mp} = - F^{\text{FV}}_{pm}$, $\widehat{F_{mp}} = -\widehat{F_{pm}}$ et $\widetilde{F_{mp}} = -\widetilde{F_{pm}}$\\
En posant $\theta_{pm} = \theta_{mp}$, on obtient : $\widetilde{{u}}_{pm}^- = \widetilde{{u}}_{mp}^+ = {u}_{mp}^{*,FV} + \theta_{mp} \frac{\Delta F_{mp}}{\gamma_{mp}}$, car $ {u}_{mp}^{*,FV} = {u}_{pm}^{*,FV}$ \\

Si l'on souhaite que la solution soit dans un ensemble convexe $\Omega$, tant que $\overline{{U}}^{n} \in \Omega$, 
il est toujours possible de trouver $\theta_{mp}$ pour que $\overline{{U}}^{n+1} \in \Omega$ car ${u}_{mp}^{*,FV}$ est combinaison convexe de $ \overline{{u}}_{mp}^{n}$ {\color{red} Voir et faire annexe}.

\subsection{choix du paramètre d'hybridation {\color{red} nom à changer}}

\paragraph{Global Maximal Principle (GMP)}
La première propriété que nous allons étudier et qui nous servira d'exemple pour la suite est comment imposer un principe du maximum global dans le cas scalaire $\Omega \subset \R^M$, $M=1$ : 
$$ \alpha \leq \overline{U}^{0} \leq \beta \implies \alpha \leq \overline{U}^{n} \leq \beta \quad \forall n$$
Pour le cas $M>1$, on regarde le cas particulier de la positivité des variables, en particulier pour le problème de la dynamique des gaz.

Si l'on souhaite obtenir une telle propriété, il nous suffit alors de trouver $\theta_{mp}$ pour avoir $\widetilde{{u}}_{mp}^{\pm,n+1}\in \left[\alpha,\beta\right]$ et comme ${u}_{mp}^{*,FV}\in \left[\alpha,\beta\right]$, on a la suite d'équivalence:
\begin{align*}
    &\widetilde{{u}}_{mp}^{\pm}\in \left[\alpha,\beta\right] \\
    &{u}_{mp}^{*,FV} \pm \theta_{mp} \frac{\Delta F}{\gamma}\in \left[\alpha,\beta\right] \\
    &\theta_{mp} \frac{\pm \Delta F}{\gamma} \in  \left[\alpha-{u}_{mp}^{*,FV},\beta-{u}_{mp}^{*,FV}\right] \\
    &\theta_{mp} \in \left[\frac{-\gamma}{|\Delta F|} \left({u}_{mp}^{*,FV} - \alpha \right), \frac{\gamma}{|\Delta F|}\left(\beta-{u}_{mp}^{*,FV}     \right) \right] \bigcap 
                     \left[\frac{-\gamma}{|\Delta F|} \left(\beta-{u}_{mp}^{*,FV}    \right), \frac{\gamma}{|\Delta F|}\left({u}_{mp}^{*,FV} - \alpha  \right) \right]
\end{align*}
En remarquant que $\frac{-\gamma}{|\Delta F|} \left({u}_{mp}^{*,FV} - \alpha \right),\frac{-\gamma}{|\Delta F|}\left(\beta-{u}_{mp}^{*,FV}\right)  \leq 0 $ et que par définition $\theta_{mp} \leq 1$.
On trouve la condition suffisante pour obtenir un schéma GMP : 
\begin{equation}
\theta_{mp} = \min \left(1, \frac{\gamma}{|\Delta F|} \min \left(\beta-{u}_{mp}^{*,FV},{u}_{mp}^{*,FV} - \alpha  \right) \right) 
\end{equation}

\paragraph{Local  Maximal Principle (LMP)}
On sait que la solution entropique d'un problème scalaire suit le principe du maximum local i.e : 
\begin{equation}
    \beta_i^n := \max_{j\in\{i-1,i+1\}} u_j^n \geq u_i^{n+1} \geq \min_{j\in\{i-1,i+1\}} u_j^n =: \alpha_i^n
\end{equation}
Il est possible d'imposer ce critère sur chaque sous-cellules d'une manière similaire qu'on impose le principe du maximum global,
la différence principale est qu'on regarde $\widetilde{{u}}_{mp}^{-}$ et $\widetilde{{u}}_{mp}^{+}$ de manière différentes car
$\widetilde{{u}}_{mp}^{-} \in \left[\alpha_m^n, \beta_m^n \right]$ tandis que $\widetilde{{u}}_{mp}^{+} \in \left[\alpha_p^n, \beta_p^n \right]$
en suivant un raisonnement similaire à la preuve du schéma GMP, et en pretant attention au signe de $\Delta F$, on obtient la condition suffisante : 
\begin{equation}
\theta_{mp} = \min \left(1, \frac{\gamma}{|\Delta F|}  \left\{
    \begin{aligned}
    \min\left( \beta_{p}^n-{u}_{mp}^{*,FV},{u}_{mp}^{*,FV} - \alpha_{m}^n \right), &\text{ si }\Delta F >0  \\
    \min\left( \beta_{m}^n-{u}_{mp}^{*,FV},{u}_{mp}^{*,FV} - \alpha_{p}^n \right), &\text{ si }\Delta F <0  
    \end{aligned}\right.
    \right)
\end{equation}

\paragraph{positivité d'un schéma sur le problème de la dynamique des gaz \\}
On donne $\Omega \subset R^3$ l'espace des phases d'un gaz comme :  
\begin{equation}
    \Omega = \{(\rho,\rho v,\rho e)\quad | \quad \rho >0, \quad \epsilon= \rho e - \rho \frac{|v|^2}{2} >0\}
\end{equation}
Là où il est facile de vérifier $\rho >0$ en reprenant la condition du maximum global. Il est moins facile d'obtenir une condition pour la positivité de $\epsilon$. 
Il n'est cependant pas impossible d'imposer un tel critère, en suivant une preuve inspiré de  \cite{VILAR2025}  \cite{Rueda-Ramírez2024} \cite{Kuzmin2020}. En supposant $\rho >0$, on cherche un critère pour obtenir $\rho \,\rho e - \frac{|\rho v|^2}{2} >0$.\\
En transposant cela dans l'écriture du schéma hybride, on cherche $\theta_{mp}$ tel que : $\widetilde{\rho}^{\pm} \widetilde{\rho e}^{\pm} - \frac{|\widetilde{\rho v}^{\pm}|^2}{2} >0$
\begin{align*}
    &\widetilde{\rho}^{\pm} \widetilde{\rho e}^{\pm} - \frac{|\widetilde{\rho v}^{\pm}|^2}{2}  \geq 0 \\
    &\left( \rho^*   \pm \theta_{mp}\frac{\Delta F^{\rho  } }{\gamma} \right) \, \left(\rho e^* \pm \theta_{mp}\frac{\Delta F^{\rho e} }{\gamma} \right) - 
     \frac{|\rho v^* \pm \theta_{mp}\frac{\Delta F^{\rho v} }{\gamma}|^2}{2}  \geq 0 \\
    &\rho^* \, \rho e^* - \frac{|\rho v|^2}{2}
    +\theta_{mp}\left(\frac{\pm\Delta F^{\rho  } }{\gamma}\rho e^*  \right)
    +\theta_{mp}\left(\frac{\pm\Delta F^{\rho e} }{\gamma}\rho^*    \right)
    -\theta_{mp}\left(\frac{\pm\Delta F^{\rho v} }{\gamma}\rho v^*    \right)
    +\theta_{mp}^2\left(\frac{\Delta F^{\rho e}\Delta F^{\rho}}{\gamma^2} \right)
    -\theta_{mp}^2\left(\frac{|\Delta F^{\rho v}|^2}{2\gamma^2} \right) \geq 0 \\
    &\rho^* \, \rho e^* - \frac{|\rho v|^2}{2} \geq 
    \mp \theta_{mp} \frac{1}{\gamma}   \left[\Delta F^{\rho} \rho e^* + \Delta F^{\rho e} \rho^* - \Delta F^{\rho v} \rho v^* \right]
    + \theta_{mp}^2 \frac{1}{2\gamma^2}\left[2\Delta F^{\rho e}\Delta F^{\rho}- |\Delta F^{\rho v}|^2\right]
\end{align*}
comme $\theta \in [0,1]$, on sait que $\theta^2 \leq \theta$ et que la condition s'écrit alors : 
\begin{equation}
\theta_{mp} = \min \left(1, \frac{\gamma}{|\Delta F|} \min \left(\beta-{u}_{mp}^{*,FV},{u}_{mp}^{*,FV} - \alpha  \right)
, \frac{M}{|B|+ min(A,0) } \right) 
\end{equation}
où $M = \rho^* \, \rho e^* - \frac{|\rho v|^2}{2}$, 
$B= \frac{1}{\gamma}   \left[\Delta F^{\rho} \rho e^* + \Delta F^{\rho e} \rho^* - \Delta F^{\rho v} \rho v^* \right]$
$A = \frac{1}{2\gamma^2}\left[2\Delta F^{\rho e}\Delta F^{\rho}- |\Delta F^{\rho v}|^2\right]$.\\

Cette condition garantie à la fois que $\widetilde{\rho}^{\pm} >0$ mais aussi que $\widetilde{\epsilon}^{\pm}>0$ 

\paragraph{Détections d'extrema locaux}
Un des problèmes de ce schéma est qu'il est très facile de se retrouver avec une convergence à l'ordre 1 
car l'on applique des critères GMP/LMP sur des zones où la solution est régulière. Il serait donc judicieux de conserver cette régularité au plus possible. \\
Il nous faut une méthode pour repérer les zones régulière de la solution. Nous allons utiliser une méthode basé sur une idée de Van Leer \cite{VANLEER1979101}, généralisée dans \cite{BURBEAU2001111} et appliquée à notre cadre dans \cite{VILAR2025}.
Le détecteur regarde la régularité de la dérivée première en regardant si sa linérairisée se trouve entre la moyenne de ses voisins. \\
\begin{wrapfigure}{r}{0.4\textwidth} 
    \usebox{\detectextrema} 
    \caption{visuel du détecteur}
\end{wrapfigure}
\begin{equation*}
    v_i(x) = \partial_x \overline{{u}_{i}} +(x-x_i) \partial_{xx} \overline{{u}_i} \qquad
\begin{aligned}
    &v_{i-1/2}^M := \max(\overline{{u}_{i-1}},\overline{{u}_{i}}) \\
    &v_{i-1/2}^m := \min(\overline{{u}_{i-1}},\overline{{u}_{i}})
\end{aligned}
\end{equation*}
\begin{align}
    &v_{i-1/2}^m \leq v_i(x_{i-1/2}) \leq v_{i-1/2}^M \\
    &v_{i+1/2}^m \leq v_i(x_{i+1/2}) \leq v_{i+1/2}^M
\end{align}
Le critère ci-dessus tient pour les cellules, mais il est plus intéressant de l'utiliser sur des sous-cellules dans notre cas.
En pratique, il est possible d'obtenir la dérivée de $u_h$ à partir de la matrice de masse et la matrice de rigidité. \\
\noindent
La première méthode consiste à comparer deux écriture de la dérivée d'un polynome dans la base polynomiale : 
$$  \partial_xf_h(x) = \sum_p^{N_k} f_h  \partial_x \sigma_p(x) = \sum_p^{N_k} f_h' \sigma_p(x) $$
Et en observant que l'on doit avoir une égalité par produit scalaire avec tout polynome de la base $\sigma_q$, on obtient une égalité de produit matrice vecteurs : 
$$ F_h K = F_h' M$$
Il est alors possible de récupérer les coordonées $F_h'$ en inversant la matrice de massse :$F_h' = M^{-1}F_h K$.\\
à partir des coordonées de la dérivée du polynome, il est possible d'obtenir sa moyenne par un passage similaire à \eqref{passage_svf}.
La deuxième méthode consiste à passer ces 2 étapes en une seul opération en calculant un opérateur de projection utilisant les dérivées des polynomes de base $\{\sigma_p'\}$: $\mathcal{P}'$ et $\mathcal{P}''$
$$\partial_x \overline{f_h}^m = \frac{1}{|S_m^i|} \int_{S_m^i}\partial_x f_h =\frac{1}{|S_m^i|}\int_{S_m^i} \sum_p^{N_k} f_p \partial_x \sigma_p(x)dx = \sum_p^{N_k}f_p \frac{1}{|S_m^i|}\int_{S_m^i} \partial_x \sigma_p(x)dx   $$
\begin{equation}
    \partial_x \overline{f_h}^m = (\mathcal{P}'F_h)_m\qquad \mathcal{P}'_{mp} =  \frac{1}{|S_m^i|}\int_{S_m^i} \partial_x \sigma_p(x)dx = \frac{\sigma_p(\tilde{x}^i_{m+1/2}) - \sigma_p(\tilde{x}^i_{m-1/2})}{|S_m^i|}
\end{equation}
\begin{equation}
    \partial_{xx} \overline{f_h}^m = (\mathcal{P}''F_h)_m\qquad \mathcal{P}''_{mp} =  \frac{1}{|S_m^i|}\int_{S_m^i} \partial_{xx} \sigma_p(x)dx = \frac{\partial_x \sigma_p(\tilde{x}^i_{m+1/2}) - \partial_x \sigma_p(\tilde{x}^i_{m-1/2})}{|S_m^i|}
\end{equation}

% \subsubsection{conservation de l'entropie}
% \paragraph{critère d'ordre 1}
% \paragraph{critère d'ordre 2}
% \paragraph{critère d'ordre arbitraire}
